Low-rank adaptation (LoRA) is widely applied to transformer models by targeting \texttt{nn.Linear} projection layers --- a recipe that transfers, practitioners assume, to any architecture with equivalent layer types.
We test this assumption on Mamba-130m, a pure state space model (SSM) whose projection layers are \texttt{nn.Linear} and whose community documentation reports 90--92\% SST-2 accuracy with projection-only LoRA.
Applying LoRA to \texttt{in\_proj}, \texttt{out\_proj}, and \texttt{x\_proj} (rank 8, 1.14\% trainable parameters), we find that SST-2 accuracy remains at 0.5092 across all three training epochs --- statistically indistinguishable from the zero-shot baseline of 0.4908 --- while MNLI accuracy degrades from 0.3463 to 0.3234, both results well below the conservative gate criterion of 0.70.
The failure is not a configuration error: LoRA installs correctly and adapter weights receive gradient updates, but those updates carry no task-discriminative information.
We attribute this pattern to the Mamba selective scan kernel, a custom CUDA operation that we interpret as acting as a gradient barrier between the classification head and the projection-layer adapters; gradient magnitudes were not directly logged, so this remains a medium-confidence mechanistic interpretation.
These findings establish that API compatibility with \texttt{nn.Linear} is not a sufficient condition for LoRA to produce learning on SSM architectures, and that Mamba parameter-efficient fine-tuning requires gradient-path-aware methods that operate inside or upstream of the SSM scan.
